\documentclass{computacion}
\titulo{Manual de Usuario --- Sistema CyAD}
\tipo{Proyecto Tecnológico}
\version{Manual de Usuario}
\trimestre{2026-Primavera}
\alumno{[TU NOMBRE]}
\matricula{[MATRÍCULA]}

\asesor{[Grado. Nombres Apellidos del Asesor]}
\categoria{[Profesor Titular/Asociado/Asistente]}
\departamento{[Departamento de Sistemas]}
\correo{[correo@azc.uam.mx]}

\coasesor{}
\categoriacoasesor{}
\departamentocoasesor{}
\correocoasesor{}

\begin{document}
\portada
\setcounter{page}{1}
\pagestyle{fancy}

\tableofcontents
\listoffigures
\newpage

\section{Introducción}
Este manual está dirigido a los tres tipos de usuario del sistema: el \textbf{profesor}, el \textbf{administrador de la División de CyAD} y el \textbf{visitante público} (cualquier persona con acceso a la URL institucional). Cada rol dispone de un \emph{dashboard} propio con las funciones que le corresponden.

Todas las capturas de este manual corresponden a la versión de producción y utilizan datos de demostración.

\section{Acceso al sistema}

Para iniciar sesión, ingrese a la URL del sistema (por ejemplo, \texttt{https://cyad.uam.mx}) y utilice el enlace ``Iniciar sesión''. La Figura~\ref{fig:m-login} muestra la pantalla de acceso.

\begin{figure}[!ht]
    \centering
    \captura[0.75\textwidth]{Imagenes/screenshots/login}
    \caption{Pantalla de inicio de sesión.}
    \label{fig:m-login}
\end{figure}

Ingrese su correo institucional (\texttt{usuario@uam.mx} o similar) y la contraseña proporcionada por el administrador. Tras iniciar sesión el sistema lo dirigirá automáticamente al \emph{dashboard} correspondiente según su rol.

\textbf{Notas de seguridad:}
\begin{itemize}
    \item Se permiten hasta 5 intentos de inicio de sesión por minuto. Si se excede, deberá esperar un minuto.
    \item La sesión expira automáticamente por inactividad; el sistema renueva el \emph{token} de acceso en segundo plano mientras haya actividad.
    \item Para restablecer una contraseña olvidada, contacte al administrador.
\end{itemize}

\section{Uso — rol Profesor}

Al iniciar sesión como profesor se muestra el \emph{dashboard} personal (Figura~\ref{fig:m-docente-dash}), con:
\begin{itemize}
    \item Tarjetas de resumen del trimestre en curso.
    \item Listas agrupadas por periodo (colapsables) de sus cartas temáticas, requisitos y autoevaluaciones.
    \item Menú lateral con acceso directo a los tres módulos.
\end{itemize}

\begin{figure}[!ht]
    \centering
    \captura[0.9\textwidth]{Imagenes/screenshots/docente_dashboard}
    \caption{Dashboard del profesor.}
    \label{fig:m-docente-dash}
\end{figure}

\subsection{Cartas temáticas}

Desde la sección ``Cartas Temáticas'' (Figura~\ref{fig:m-cartas-lista}) puede crear, editar, publicar y consultar la vista pública de cada carta.

\begin{figure}[!ht]
    \centering
    \captura[0.9\textwidth]{Imagenes/screenshots/docente_cartas_lista}
    \caption{Lista de cartas temáticas del profesor.}
    \label{fig:m-cartas-lista}
\end{figure}

\subsubsection{Crear una carta temática}
\begin{enumerate}
    \item Haga clic en \textbf{``Nueva carta temática''}.
    \item Busque la UEA por clave o nombre; seleccione una opción del autocompletado.
    \item El sistema asigna automáticamente el periodo activo para cartas.
    \item Complete los campos del grupo (nombre, ID, horario, modalidad).
    \item Complete el contenido: descripción, objetivos, contenidos, evaluación, bibliografía. Todos son texto libre; los que deje vacíos no aparecerán en la vista pública.
    \item Presione \textbf{``Guardar''} para conservar como \emph{borrador} o \textbf{``Publicar''} para que sea visible en la vista pública.
\end{enumerate}

\begin{figure}[!ht]
    \centering
    \captura[0.9\textwidth]{Imagenes/screenshots/docente_carta_form}
    \caption{Formulario de creación/edición de carta temática.}
    \label{fig:m-carta-form}
\end{figure}

Una carta publicada puede ser \emph{despublicada} y editada nuevamente; el historial se conserva en la base de datos.

\subsection{Requisitos de recuperación}

El flujo es análogo al de cartas temáticas. Desde ``Requisitos de Rec.'' \textbf{→ ``Nuevo requisito''} se abre el formulario correspondiente (Figura~\ref{fig:m-req-form}).

\begin{figure}[!ht]
    \centering
    \captura[0.9\textwidth]{Imagenes/screenshots/docente_requisito_form}
    \caption{Formulario de requisitos de recuperación.}
    \label{fig:m-req-form}
\end{figure}

\subsection{Autoevaluación docente}

Cuando el administrador abre un periodo de autoevaluación, aparece en la sección ``Autoevaluación'' el formulario correspondiente (Figura~\ref{fig:m-autoeval}). El formulario admite ocho tipos de pregunta:

\begin{itemize}
    \item Texto corto, texto largo.
    \item Opción única (radios), casillas (múltiple).
    \item Sí/No, escala lineal.
    \item Lista desplegable, cuadrícula.
\end{itemize}

\begin{figure}[!ht]
    \centering
    \captura[0.9\textwidth]{Imagenes/screenshots/docente_autoevaluacion}
    \caption{Formulario de autoevaluación con renderizado dinámico.}
    \label{fig:m-autoeval}
\end{figure}

Puede guardar respuestas parciales (\emph{en progreso}) y regresar más tarde. Al terminar, presione \textbf{``Enviar''}; el sistema calcula automáticamente el puntaje ponderado y le muestra el nivel de desempeño alcanzado. Una vez enviada, la autoevaluación no puede modificarse.

\section{Uso — rol Administrador}

El \emph{dashboard} del administrador (Figura~\ref{fig:m-admin-dash}) muestra indicadores globales: total de profesores, cumplimiento por departamento, autoevaluaciones enviadas, etc.

\begin{figure}[!ht]
    \centering
    \captura[0.9\textwidth]{Imagenes/screenshots/admin_dashboard}
    \caption{Dashboard del administrador.}
    \label{fig:m-admin-dash}
\end{figure}

\subsection{Gestión de profesores}
Desde ``Profesores'' (Figura~\ref{fig:m-admin-prof}) puede:
\begin{itemize}
    \item Crear un profesor con correo, nombre y departamento.
    \item Editar los datos del perfil.
    \item Restablecer la contraseña (se genera una temporal que debe entregarle al profesor).
    \item Desactivar cuentas.
\end{itemize}

\begin{figure}[!ht]
    \centering
    \captura[0.9\textwidth]{Imagenes/screenshots/admin_profesores}
    \caption{Gestión de profesores.}
    \label{fig:m-admin-prof}
\end{figure}

\subsection{Constructor de formularios de autoevaluación}

El constructor (Figura~\ref{fig:m-admin-builder}) permite crear formularios versionados. Cada formulario se compone de \emph{secciones}, cada sección de \emph{preguntas}, y algunas preguntas de \emph{opciones}.

\begin{figure}[!ht]
    \centering
    \captura[0.9\textwidth]{Imagenes/screenshots/admin_formulario_builder}
    \caption{Constructor de formularios (form-builder).}
    \label{fig:m-admin-builder}
\end{figure}

\textbf{Reglas de validación:} la suma de los pesos de las secciones debe ser exactamente 100~\%. El sistema no permite publicar un formulario que no cumpla esta regla.

\subsection{Reportes}

Desde ``Reportes'' (Figura~\ref{fig:m-admin-rep}) puede consultar cumplimiento por departamento, licenciatura, profesor o periodo, y descargar cada tabla en formato XLSX (Excel).

\begin{figure}[!ht]
    \centering
    \captura[0.9\textwidth]{Imagenes/screenshots/admin_reportes}
    \caption{Panel de reportes.}
    \label{fig:m-admin-rep}
\end{figure}

\subsection{Catálogos institucionales}
Desde el menú ``Catálogos'' se administran: departamentos, licenciaturas, posgrados, áreas, UEAs y periodos. Cada periodo tiene banderas por recurso (cartas, requisitos, autoevaluación) que determinan qué se puede capturar en ese trimestre.

\section{Uso — Vista pública}

La vista pública ``Explorar'' (Figura~\ref{fig:m-publico}) está disponible sin iniciar sesión. Presenta una cascada:

\begin{enumerate}
    \item Selección de periodo (trimestre).
    \item Selección de programa (licenciatura o posgrado).
    \item Selección de UEA.
    \item Listado de cartas y requisitos publicados para esa UEA en ese periodo.
\end{enumerate}

\begin{figure}[!ht]
    \centering
    \captura[0.9\textwidth]{Imagenes/screenshots/publico_explorar}
    \caption{Vista pública ``Explorar''.}
    \label{fig:m-publico}
\end{figure}

Al hacer clic en un documento, se abre en una nueva pestaña con formato imprimible.

\section{Preguntas frecuentes}

\begin{description}
    \item[¿Qué pasa si publico una carta con errores?] Puede despublicarla, corregirla y volver a publicarla; la URL pública se mantiene.
    \item[¿Puedo tener varias cartas para la misma UEA en el mismo trimestre?] Sí, si son grupos distintos (identificados por \emph{Nombre del grupo} + \emph{ID del grupo}).
    \item[¿Se puede reabrir una autoevaluación ya enviada?] No desde la interfaz. El administrador puede reabrirla desde la base de datos si es estrictamente necesario.
    \item[¿Cómo se calcula el puntaje?] Cada pregunta tiene un peso; el peso de la sección multiplica al de sus preguntas; la suma final da el porcentaje total. Los niveles de desempeño (por ejemplo, \emph{Insuficiente}, \emph{Suficiente}, \emph{Satisfactorio}, \emph{Sobresaliente}) se definen por rangos configurables por el administrador.
    \item[¿La vista pública muestra los datos personales del profesor?] No. Muestra el nombre público (que el administrador puede editar), la UEA y el contenido académico. Datos como correo o matrícula quedan privados.
\end{description}

\end{document}
